Table~\ref{tab:pathtable} uses the common validation protocol, including update zero and every 100th continuation update. Each selected model satisfies the flow threshold; Appendix~\ref{app:pairs-v102} gives full objectives, denominators, the principal four-method comparison, and the earlier pointwise explorations. The test set was expanded from 24 to 96 demand paths --- twenty-four paths per active-query group, generated by the same rule under a fresh seed --- with the trained models, training data and validation selection frozen, and no result on the expanded set used for model selection or method modification; paths are the independent unit, and active queries and seeds are averaged within a path before any summary. The original 24-path evaluation was complete before the expansion was generated, so the 96-path comparison is an expanded robustness evaluation under a frozen protocol rather than a preregistered confirmatory test.
\begin{table}[t]\centering\footnotesize
\caption{Prediction and spatial explanation on the same 96 test demand
paths (24 per group). Entries are group medians of path-level values;
flow errors are additionally averaged over the 41 states within each
path. Geometric: the maximum account error over the geometric family,
as a percentage of the reference geometric-family maximum.
Unrestricted and bound: the exact maximum $\mathcal K_j(e_j)$ over all
separating regions and the integrated bound of
Eq.~\eqref{eq:localboundarybound}, as percentages of the reference
column's boundary seminorm. The geometric and unrestricted columns use
different denominators and are not one inequality; the bound chain
$\mathcal K_j(e_j)\le B_j$ holds in absolute units for the shared
quadrature (text below). Improved:
paths on which response supervision lowers the error.}\label{tab:pathtable}
\setlength{\tabcolsep}{3.4pt}
\begin{tabular*}{\tblwidth}{@{\extracolsep{\fill}}rcccccc cc cc cc cc}
\toprule
& \multicolumn{2}{c}{Flow pred.\ (\%)} & \multicolumn{2}{c}{Column (\%)} & \multicolumn{2}{c}{Geometric (\%)} & \multicolumn{2}{c}{Unrestricted (\%)} & \multicolumn{2}{c}{Bound (\%)} & \multicolumn{2}{c}{Improved}\\
\cmidrule(lr){2-3}\cmidrule(lr){4-5}\cmidrule(lr){6-7}\cmidrule(lr){8-9}\cmidrule(lr){10-11}\cmidrule(lr){12-13}
Active & flow & resp. & flow & resp. & flow & resp. & flow & resp. & flow & resp. & col. & geo.\\
\midrule
1 & 0.032 & 0.029 & 8.13 & 4.92 & 2.95 & 1.78 & 4.54 & 2.95 & 4.62 & 3.00 & 24/24 & 23/24\\
2 & 0.032 & 0.029 & 21.09 & 15.67 & 8.62 & 6.19 & 15.48 & 12.65 & 15.67 & 12.80 & 22/24 & 21/24\\
3 & 0.033 & 0.029 & 21.53 & 15.46 & 9.07 & 6.01 & 16.36 & 12.39 & 16.70 & 12.81 & 22/24 & 22/24\\
6 & 0.033 & 0.030 & 22.45 & 14.22 & 9.27 & 5.77 & 19.43 & 12.28 & 20.61 & 13.48 & 24/24 & 24/24\\
\midrule
All 96 & 0.032 & 0.029 & 21.18 & 14.19 & 8.74 & 5.75 & 15.65 & 12.08 & 16.34 & 12.83 & 92/96 & 90/96\\
\bottomrule\end{tabular*}\end{table}
On the 96 test demand paths, response supervision lowers the median
attribution-column error from 21.18\% to 14.19\% (a relative reduction
of one third) and the geometric-boundary error from 8.74\% to 5.75\%,
while both models retain flow-prediction errors near $0.03\%$ (medians
$0.0322\%$ and $0.0292\%$). The improvement holds in every
active-query group. Path-level paired differences quantify the
uncertainty: the median paired column-error difference is $-5.03$
percentage points with a path-level bootstrap $95\%$ interval of
$[-6.01,-4.48]$ (10,000 resamples), improving 92 of 96 paths; the
paired geometric difference has median $-2.06$ with interval
$[-2.76,-1.65]$, improving 90 of 96, with every group's interval
excluding zero and leave-one-path-out medians varying within
$[-5.11,-4.94]$ and $[-2.11,-2.00]$
(Appendix~\ref{supp:learningdetail} reports the group-level intervals).
The system-contribution difference is secondary and opposite in sign
(median $+0.45$, interval $[0.01,0.74]$; both methods below $3.5\%$),
confirming that scalar completeness and spatial explanation accuracy
measure different properties. On the original 24-path subset, the
same-path decomposition in Eq.~\eqref{eq:coldecomp} attributes most of
the attribution discrepancy to demand-response error rather than
marginal-cost weight error, which is why the unweighted Jacobian loss
is a workable proxy for the weighted quantity of
Eq.~\eqref{eq:localboundarybound}: for both models the response
component has the larger norm on all 24 paths, with path-median
response and weight norms 256.1 and 4.1 under flow supervision and
158.1 and 3.8 under response supervision.

The bound chain closes directly on the test paths: computed with the
same $41$-state trapezoidal quadrature for both sides, the exact
unrestricted errors $\mathcal K_j(e_j)$ lie below the integrated
bounds $B_j$ of Eq.~\eqref{eq:localboundarybound} at machine precision
(maximum relative chain violation $1.6\times10^{-15}$ over every
path, query, seed and method; this verifies the inequality for the
shared quadrature, and quadrature accuracy is reported separately in
Appendix~\ref{supp:accuracy}), and the chain is tight --- the median
ratio $\mathcal K_j(e_j)/B_j$ is $0.95$ for flow supervision (quartiles
$0.91$--$0.98$) and $0.95$ for response supervision (quartiles
$0.90$--$0.97$). The theory's integrated bound is therefore an
accurate predictor of the realized worst boundary error, and response
supervision lowers the geometric, unrestricted and bound levels in
parallel. The geometric family retains about half of the unrestricted
error (median ratio $E^{\mathrm{geo}}/\mathcal K_j(e_j)$ $0.51$ and
$0.48$), so the single-family errors reported throughout should be
read as lower bounds on the worst separating-region error.
Among the evaluated supervision objectives, response supervision
produced the most consistent gains across active-query groups; the
endpoint-edge objectives give no overall improvement and select the
common initial model for three of five seeds
(Appendix~\ref{app:pairs-v102}). Sign reliability on the separate
frozen six-query baseline segment is boundary specific and is retained
in Appendix~\ref{supp:screenlinelearned}.
